\section{Mapa del sistema, límites señalados en la Q\&A y agenda de investigación}

Una vez concluidos los cuatro casos, la clase los vuelve a ubicar dentro de un solo sistema robótico y, en la sesión de preguntas, expone sus límites por iniciativa propia: la skill library sigue siendo pequeña, la imitation es solo un punto de partida, el compositional scaling no tiene un umbral y un contexto visual de seis cuadros ya está cerca de la capacidad del modelo.

\subsection{What's Next: primero, regresar al nivel del sistema}

El último punto de la Agenda no presenta otro artículo, sino que combina los local successes. Cualquier afirmación sobre un «foundation model para robots» debe responder quién se encarga de skill learning, planning, low-level control, data augmentation y object representation, cómo se transmiten los datos entre las interfaces y cómo se retroalimentan las fallas.

\lecturefigure{slide-42-agenda-whats-next.jpg}{Agenda: de los cuatro trabajos de regreso al sistema robótico completo}{Diapositiva V042 recuperada del video oficial; Stanford CS25 Lecture 15}

\subsection{Component Map: dónde se ubica cada Foundation Model}

La tabla del sistema coloca RT-1 en skill learning; SayCan / Inner Monologue, en planning; Code as Policies, en low-level control; DIAL, en data augmentation, y NLMap, en object-centric representations. La recipe de la parte superior subraya que todos estos métodos dependen de large diverse offline data, de una high-capacity architecture y de language glue.

\lecturefigure{slide-43-system-component-map.jpg}{Component map del sistema: RT-1, SayCan, Inner Monologue, Code as Policies, DIAL y NLMap}{Diapositiva V043 recuperada del video oficial; Stanford CS25 Lecture 15}

\begin{importantbox}{Language es una API, no una inteligencia monolítica}
El lenguaje permite que skills, planes, scene descriptions y labels compartan una interfaz legible, por lo que es fácil conectar modelos externos; sin embargo, cada módulo sigue necesitando sus propios data, evaluation y failure handling. Fusionar todos los módulos en un solo decoder puede simplificar la interfaz, pero no elimina por sí solo los problemas de grounding ni de control en tiempo real.
\end{importantbox}

\subsection{2016--2023: cómo cambió la pregunta}

La timeline escribe bajo el eje horizontal las preguntas de tres etapas: en 2016--2021, cómo aprende un robot real de extremo a extremo; en 2021--2022, cómo escalar a muchas tasks; en 2022--2023, cómo aprovechar los foundation models para acelerar la robotics. SayCan, Inner Monologue, DIAL, Code as Policies y NLMap, enumerados a la derecha, son componentes de un sistema, no un end-to-end model unificado.

\lecturefigure{slide-44-robotics-progress-timeline.jpg}{Avances en la investigación robótica: de end-to-end y multi-task scaling al foundation-model leverage}{Diapositiva V044 recuperada del video oficial; Stanford CS25 Lecture 15}

\subsection{Closing y proyectos públicos}

La línea de tiempo ya devolvió RT-1, SayCan, Inner Monologue y DIAL a los problemas que cada uno resuelve; el cierre, a su vez, vuelve a anclar estos sistemas en artículos, código y demostraciones que pueden revisarse. Esta sección no añade otro método, sino que revisa qué puntos de entrada verificables deja la clase completa. Al leer la última diapositiva conviene hacer corresponder cada proyecto público con los módulos de data, planning, feedback y grounding vistos antes, y a la vez evitar tomar una demo cuidadosamente seleccionada como prueba completa de operación autónoma a largo plazo.

La Thank-you slide conserva los puntos de entrada de cuatro proyectos y un ejemplo de un robot que limpia un refresco de cola. Recuerda al lector que la evidencia de esta clase proviene de artículos, código y demos consultables, no de una visión abstracta. En la demostración, la instrucción larga se sigue descomponiendo en una secuencia finita de habilidades, lo que muestra el límite entre language planning y la skill library.

\lecturefigure{slide-45-thank-you.jpg}{Cierre: proyectos SayCan, Inner Monologue, RT-1 y DIAL, y demostración de la tarea de limpieza}{Diapositiva V045 recuperada del video oficial; Stanford CS25 Lecture 15}

\subsection{Q\&A uno: la Generalization no tiene un umbral único}

Un estudiante pregunta qué escala requiere la generalización composicional. El ponente responde que depende de la task definition, de la dataset scale y del concepto; hoy no hay una cifra que prediga la emergence a través de colores, objetos y habilidades. Su conjetura personal es que tal vez falten uno o dos órdenes de magnitud de datos, pero aclara de forma explícita que se trata de un juicio subjetivo.

\begin{warningbox}{No conviertas «uno o dos órdenes de magnitud» en una Scaling Law}
Esta frase proviene de una rough intuition expresada durante las preguntas, sin experimentos controlados, sin definición del eje horizontal y sin intervalos de confianza. Sirve para expresar que «la escala actual sigue siendo pequeña», pero no para predecir que con cierto número de parámetros o cierto episode count surgirán necesariamente capacidades robóticas generales.
\end{warningbox}

\subsection{Q\&A dos: el cuello de botella de SayCan son las Grounded Skills}

Las affordance/value functions de SayCan solo pueden asignar puntuaciones a una skill library finita. Si el robot tiene únicamente tres habilidades, el planner solo puede producir combinaciones de esas tres; agrandar el modelo de lenguaje no agrega nuevas primitive de acción. Por ello, la instruction coverage de alto nivel depende del número, la precisión y la cobertura de estados de las habilidades de bajo nivel.

\teachervoice{El ponente lo llama el «100\% bottleneck» de SayCan: cuántos platillos hay en el buffet del planner determina cuántos objetivos puede combinar. Esta comparación es más precisa que decir «el LLM se encarga del planning», porque devuelve el límite de capacidad a los datos del robot y al controller.}

\subsection{Q\&A tres: la Imitation es una Existence Proof, no el final}

En su momento, RT-1 usaba sobre todo imitation loss. Como investigador de RL, el ponente quisiera incorporar offline improvement o RL, pero se eligió la multi-task imitation porque es estable, escalable y ya demostró que funciona con datos reales a gran escala. Desde el punto de vista de ingeniería, establecer primero una base policy confiable y después estudiar una optimization más compleja es más verificable que intentar resolver todos los problemas al mismo tiempo.

\begin{knowledgebox}{Por qué no usar Online RL directamente}
En robots reales, la reward es escasa, las fallas son costosas y la paralelización de datos está limitada por el hardware. La offline imitation renuncia a la exploración activa a cambio de un entrenamiento estable y de un dataset reutilizable. Si en el futuro se incorpora RL, todavía habrá que resolver el distribution shift, las restricciones de seguridad y la extrapolación de value fuera de línea.
\end{knowledgebox}

\subsection{Q\&A cuatro: seis cuadros ya están cerca del Context Limit}

Las tareas largas requieren memoria del orden de minutos, pero los tokens de imágenes de alta dimensión agotan pronto el context. Los seis cuadros de RT-1 solo cubren el movimiento de corto plazo y no pueden almacenar una trayectoria completa ni una few-shot demonstration. El ponente propone el retrieval transformer u otros memory mechanisms como posibles direcciones, y afirma de forma explícita que hoy no es realista meter trayectorias completas en el context.

\teachervoice{La última limitación es muy concreta: no se trata de que «el modelo aún no sea lo bastante inteligente», sino de que seis imágenes ya se acercan al context budget. Las tareas corporizadas de larga duración requieren compresión, recuperación, resúmenes de estado o memory externa; ampliar sin más el prompt eleva a la vez la latency y el riesgo para el control en tiempo real.}

\subsection{Resumen de la sección}

El mapa del sistema muestra que los cuatro trabajos cubren, respectivamente, las interfaces de skill, plan, feedback y data; las preguntas, en cambio, señalan que no se combinan por sí solos en un robot de propósito general. Los verdaderos cuellos de botella pendientes son la grounded skill coverage, los algoritmos de aprendizaje, una generalization predecible y la memory de largo plazo.
